\section{Experiments}

\subsection{Settings}

\subsubsection{Tasks}
We conduct experiments on both standard LLM benchmarks and more complex agentic tasks. We describe the tasks, datasets, and evaluation metrics as follows:

\begin{table}[htb]
  \centering
      \caption{\textbf{Results on Standard LLM Benchmarks.}}
      \label{tab:benchmark_results}
\renewcommand\tabcolsep{8pt}
\renewcommand\arraystretch{1.2}
  \begin{tabular}{l|cc|cc|cc} %
\Xhline{1.5pt}
\textbf{Methods} & \multicolumn{2}{c|}{\textbf{HotPotQA}} & \multicolumn{2}{c|}{\textbf{MATH}} & \multicolumn{2}{c}{\textbf{HumanEval}} \\
                 & \textbf{GPT-3.5}  & \textbf{GPT-4} & \textbf{GPT-3.5}  & \textbf{GPT-4} & \textbf{GPT-3.5}  & \textbf{GPT-4} \\
\Xhline{1.5pt}
\bf GPTs &   24 / 38.8 & 33 / 44.3  & 23.2 & 53.1 & 59.2  &  71.7 \\
\bf Agents  & 27 / 37.5 & 39 / 49.8  & 23.8 & 56.0 & 59.5 & 85.0 \\
\bf Agents w/ AutoPE & 29 / 39.8 & 38 / 50.3  & 22.5  & 57.2  & 63.5  &  82.3 \\
\bf DSPy & 35 / 43.9 & 40 / 50.5  & 17.3 & 48.4 & \bf 66.7  & 77.3 \\ \hline

\bf Ours  & \bf 35 / 44.8 & \bf 41 / 54.0  & \bf 38.8  & \bf 60.7  & 64.5  & \bf 85.8   \\ \hline
\Xhline{1.5pt}
\end{tabular}
\end{table}
\vspace{-10pt}



\paragraph{Standard Benchmarks}
We conduct experiments on standard benchmarks for LLMs including HotpotQA~\citep{yang-etal-2018-hotpotqa}, MATH~\citep{hendrycksmath2021}, and HumanEval~\citep{chen2021codex}. HotPotQA is a multi-hop QA task challenging for rich background knowledge. We use the ``hard'' split in the dataset since we find it to be more challenging for language agents. MATH is a collection of challenging competition mathematics problems. HumanEval is an evaluation set that requires LLMs or agents to synthesize programs from docstrings. As for evaluation metrics, we use F1 and exact match for HotPotQA, accuracy for MATH, and Pass@1 for HumanEval. Tools are disabled in these datasets to ensure the results comparison is meaningful with existing literature on these tasks. 

\paragraph{Complex Agent Tasks}
We consider \textbf{creative writing} and \textbf{software development} as two complex agentic tasks. For the creative writing task, we follow~\citet{yao2023tree} and give 4 random sentences to the agents and ask them to write a coherent passage with 4 paragraphs that end in the 4 input sentences respectively. Such a task is
open-ended and exploratory, and challenges creative thinking as well as high-level planning. We use GPT-4 score to evaluate the passages following~\citep{yao2023tree}. The software development task, on the other hand, requires the agent system to develop an \textit{executable} software given a simple product requirement document (PRD). We evaluate the compared agents according to the \textit{executability} of the generated software, which is quantified by numerical scores ranging from $1$ to $4$, corresponding to increasing levels of execution capability. Specifically, a score of $1$ signifies execution failure, $2$ denotes successful code execution, $3$ represents conformance to the anticipated workflow, and $4$ indicates flawless alignment with expectations.




%\paragraph{Poker}
% We adopted the rlcard library to construct a highly simulated Texas poker environment, inspired by the research framework of \citep{zhang2024agent}. However, for the convenience and efficiency of the experiment, we selected a simplified version of unlimited bet Texas poker and simplified the actions that players can take into four basic strategies: raise, call, check, and fold, in order to focus on evaluating and comparing core strategies.In order to ensure the fairness and comprehensiveness of the experiment, the entire experimental design includes 150 rounds of matches, which are divided into three stages, each with 50 rounds. During this stage, we systematically exchanged the positions of the participating agents (Agents 1.0, Agents 2.0, and GPTs) on the playing table. This arrangement ensures that all agents have gone through exactly the same game sequence, thus thoroughly testing the stability and adaptability of their strategies.Finally, we rely on the total \textit{revenue} of each agent in 150 rounds of games as our evaluation result.(see Table ~\ref{tab:Texas Poker})

% \begin{table}[htbp]
% \centering
% \caption{The Result of Texas Poker}
% \label{tab:Texas Poker}
% \begin{tabular}{|c|c|c|c|}
% \hline
% \textbf{Agents} & \textbf{Agents1.0} & \textbf{Agents2.0} & \textbf{GPTs} \\ \hline
% Revenue &  &  &  \\ \hline
% \end{tabular}
% \end{table}



\subsubsection{Baselines}

We compare our proposed method against the following baselines:
\begin{itemize}
    \item \textbf{GPTs}: a simple baseline that use GPT and a carefully designed prompt; 
    \item \textbf{Agents}: a language agent method implemented using the Agents~\citep{zhou2023agents} framework\footnote{We have tested with other agent frameworks such as OpenAgents and AgentVerse and got similar results.} with a carefully designed prompts, tools, and pipeline;
    \item \textbf{DSpy}: a LLM pipeline optimization framework that can search the best combination of prompt components. It is not applicable for complex agent tasks where the evaluation metric can not be defined in equation and code;
    \item  \textbf{Agents + AutoPE}: a variant where the prompt in each node of the agent pipeline is optimized by an LLM following the method described in~\citet{yang2024large}. It does not involve language gradient back-propagation and language gradient-based optimization.
\end{itemize}

We report results with both GPT-3.5 and GPT-4. We use the \texttt{gpt-3.5-turbo-0125} endpoint for GPT-3.5 and the \texttt{gpt-4-turbo-0409} endpoint for GPT-4. As for our approach, we start with the \textbf{Agents} baseline and then conduct \baby on top of it.

% We compare our proposed method against the following baselines: (1) \textbf{GPTs}: a simple baseline that use GPT and a carefully designed prompt; (2) \textbf{Agents}: a language agent method implemented using the Agents~\citep{zhou2023agents} framework\footnote{We have tested with other agent frameworks such as OpenAgents and AgentVerse and got similar results.} with a carefully designed prompts, tools, and pipeline; (3) \textbf{DSpy}: a LLM pipeline optimization framework that can search the best combination of prompt components. It is not applicable for complex agent tasks where the evaluation metric can not be defined in equation and code; (4) \textbf{Agents + AutoPE}: a variant where the prompt in each node of the agent pipeline is optimized by an LLM following the method described in~\citet{yang2024large}.

% \paragraph{Agents + AutoPE}

\subsection{Results}


\begin{table}[htb]
  \centering
  \begin{minipage}{.45\linewidth}
      \caption{Results on software development.}
      \label{tab:software}
      \begin{tabular}{l|ccc} %
      \Xhline{1.5pt}
      \textbf{Task} & \textbf{GPTs} & \textbf{Agents} & \textbf{Ours} \\
      \hline
      \bf Flappy bird & 2 & 2 & 3 \\
      \bf Tank battle game & 1 & 2 & 4 \\
      \bf 2048 game & 1 & 2 & 4 \\
      \bf Snake game & 2 & 3 & 4 \\
      \bf Brick breaker game & 2 & 3 & 4\\ \hline
      \bf Average score & 1.6 & 2.4 & \bf 3.8 \\
      \Xhline{1.5pt}
      \end{tabular}
  \end{minipage}%
  \hfill
  \begin{minipage}{.45\linewidth}
      \caption{Results on creative writing.}
      \label{tab:writing}
      \begin{tabular}{l|cc} %
      \Xhline{1.5pt}
      \textbf{Methods} & \textbf{GPT-3.5} & \textbf{GPT-4} \\
      \hline
      \bf GPTs & 4.0 & 6.0 \\
      \bf Agents & 4.2 & 6.0 \\
      \bf Agents w/ AutoPE & 4.4 & 6.5 \\
      \bf ToT & 3.8 & 6.8 \\ \hline
      \bf Ours & \bf 6.9 &  \bf 7.4 \\
      \Xhline{1.5pt}
      \end{tabular}
  \end{minipage}
\end{table}
  \vspace{-15pt}


\paragraph{Results on LLM Benchmarks}
The results on standard LLM benchmarks are shown in Table \ref{tab:benchmark_results}. We can see that the proposed \baby framework consistently improves over all compared methods. The performance improvement on MATH, a competition-level benchmark, is especially large. In contrast, conventional LLM-based prompt optimization method (Agents w/ AutoPE) and the search-based prompt optimization approach (DSPy) are not as stable: they results in good performance improvements in some cases but leads to significant performance degration in some other cases. This suggests that the \baby framework is more robust and can optimize the overall performance of language agents more effectively.


\paragraph{Results on Complex Tasks}
We present the results on software development and creative writing in Table \ref{tab:software} \& \ref{tab:writing}, respectively. We can see that our approach significantly outperforms all compared baselines on both tasks with a even larger performance gap compared to that on conventional LLM benchmarks. Interestingly, our approach even outperforms tree-of-throught, a carefully designed prompt engineering and inference algorithm, on the creative writing task. We find that our approach successfully finds the plan, write, and revision pipeline and the prompts are very well optimized in each step. We also find that the \baby framework recovers similar standard operation procedure developed in MetaGPT~\citep{hong2023metagpt}, an agent framework specifically designed for software development. This confirms the effectiveness of the proposed \baby framework on real-world tasks where there is no ground truth and the overall performance cannot be calculated by equations or codes, as contrary to search-based algorithms such as DSPy.

\subsection{Case Study \& Analysis}

\begin{figure}[tp]
    \centering
    \includegraphics[width=0.75\linewidth]{figs/case.pdf}
    \caption{An case study conducted on creative writing task.}
    \label{fig:case}
    \vspace{-15pt}
\end{figure}

We then show a case study for the optimization dynamics of the \baby framework in Figure \ref{fig:case}. We can see that our approach can effectively do prompt engineering and designing the agent pipeline in the way a human expert develops language agents.

Moreover, we find that the initialization of the agent system has non-negligible impacts on the final performance, just as the initialization of a neural nets is important for training. In general, we find that it is generally helpful to initialize the agent in the simplest way and let the symbolic optimizers to do the optimization. In contrast, the performance tends to become unstable if the initial agent system is over-engineered. A natural extension of this observation is that maybe we can do some kind of pre-training on large-scale and diverse tasks as a versatile initialization for general-purpose agents and then adapt it to specialized tasks with \baby. We also find that the success of our approach is more significant and stable on complex real-world tasks compared to that on standard benchmarks where the performance is evaluated by accuracy or F1. This suggests that future research on agent learning should focus more on real-world tasks, and the agent research community should work on building a benchmark focusing on agent learning evaluation that consists of diverse complex agentic tasks and investigating robust approaches to measure progress.